
\medskip
\phantomsection\label{par:sga3-VIA-3.3}%
\noindent\textbf{3.3.---}%
\footnote{\label{note:sga3-VIA-42}Editor's note: we have added this
subsection.}
In \hyperref[thm:sga3-VIA-3.2]{Theorem~3.2}, the hypothesis that \(G\) be flat can be omitted when
\(u:F\to G\) is a monomorphism. This generalization is mentioned in
Expos\'e~VI\textsubscript{B}, Remark~9.3(b), and also in
\textup{[Ray67a]}, Example~a)\SGARefLink{thm:sga3-VIA-3.2:item:i}{(i)}, p.~82. The proof, given in
\textup{[An73]}, \SGARefLink{sec:sga3-VIA-4}{Theorem~4}, follows from \hyperref[thm:sga3-VIA-3.2]{Theorem~3.2} and from the
following theorem of Grothendieck (mentioned in \textup{[Ray67a]},
Theorem~1(ii), and proved in \textup{[DG70]}, \S~III.2, \SGARefLink{prop:sga3-VIB-7.1}{7.1}).
If \(X\) is a scheme and \(R\) an equivalence relation on \(X\), we
write \(\widetilde{X/R}\) for the \((\mathrm{fppf})\) quotient sheaf
of \(X\) by \(R\) (cf.\ \SGARefLink{sga3:iv:proposition:4:4:9}{Expos\'e~IV, 4.4.9}).

\medskip
\phantomsection\label{thm:sga3-VIA-3.3.1}%
\noindent\textbf{Theorem 3.3.1 (Grothendieck).}
\textit{Let \(A\) be a ring, let \(X\) be an \(A\)-scheme, and let
\[
  d_0,d_1:R\rightrightarrows X
\]
be an \(A\)-equivalence relation on \(X\), with each \(d_i\)
faithfully flat and of finite presentation. Let \(X_0\) be a saturated
closed subscheme of \(X\), defined by a nilpotent ideal, and let
\(R_0\) be the equivalence relation induced by \(R\) on \(X_0\).
If the \((\mathrm{fppf})\) quotient sheaf
\(\widetilde{X_0/R_0}\) is representable by an \(A\)-scheme, then so
is \(\widetilde{X/R}\).}

For the proof, see \textup{[DG70]}, \S~III.2, \SGARefLink{prop:sga3-VIB-7.1}{7.1}. Let us now return
to the case where \(A\) is an Artinian local ring. Let
\(u:F\to G\) be a quasi-compact morphism between \(A\)-groups locally
of finite type, and suppose in addition that \(u\) is a monomorphism.
Then \(u\) is a closed immersion by \SGARefLink{prop:sga3-VIA-2.5.2}{2.5.2}.

We may now state the following variant of \hyperref[thm:sga3-VIA-3.2]{Theorem~3.2}.

\medskip
\phantomsection\label{thm:sga3-VIA-3.3.2}%
\noindent\textbf{Theorem 3.3.2.}
\textit{Let \(A\) be an Artinian local ring, let \(G\) be an
\(A\)-group locally of finite type, and let \(F\) be a closed
subgroup of \(G\), flat over \(A\).}%
\footnote{\label{note:sga3-VIA-43}Editor's note: for an example in
which \(F\) is not flat and \(\widetilde{F\backslash G}\) is not
representable, see \textup{[DG70]}, \S~III.3, no.~\SGARefLink{prop:sga3-VIB-3.3}{3.3}.}
\textit{Then:}
\begin{enumerate}
\item[\SGAInlineRefTarget{thm:sga3-VIA-3.3.2:item:i}(i)] \textit{The \((\mathrm{fppf})\) quotient sheaf
\(\widetilde{F\backslash G}\) is representable by a separated
\(A\)-scheme \(F\backslash G\), locally of finite type; moreover, the
sequence}
\begin{center}
  \phantomsection\label{diag:sga3-VIA-3.3.2-exact-sequence}%
  \begin{tikzcd}[column sep=large]
    F\times_A G
      \arrow[r,shift left=1.2ex]
      \arrow[r,shift right=1.2ex]
    &
    G
      \arrow[r,"p"]
    &
    F\backslash G
  \end{tikzcd}
\end{center}
\textit{is exact in the category of all ringed spaces.}

\item[\SGAInlineRefTarget{thm:sga3-VIA-3.3.2:item:ii}(ii)] \textit{The morphism
\[
  F\times_A G
    \xrightarrow{(\lambda,\operatorname{pr}_2)}
  G\times_{(F\backslash G)}G
\]
is an isomorphism, and
\(p:G\to F\backslash G\) is faithfully flat and locally of finite
presentation; hence \(p\) is an \(F\)-torsor locally trivial for the
\((\mathrm{fppf})\) topology.}

\item[\SGAInlineRefTarget{thm:sga3-VIA-3.3.2:item:iii}(iii)] \textit{If \(G\) is flat (respectively, of finite type;
respectively, smooth) over \(A\), then so is \(F\backslash G\).}

\item[(iv)]\phantomsection\label{concl:sga3-VIA-3.3.2-iv} \textit{\(X=F\backslash G\) is equipped with a right
action of \(G\) such that \(p(e)\cdot g=p(g)\) for every \(g\in G\).
Consequently, the connected components of \(X\) are of finite type
over \(A\), irreducible, and all of dimension
\(\dim G-\dim F\).}

\item[(v)]\phantomsection\label{concl:sga3-VIA-3.3.2-v} \textit{If, in addition, \(F\) is a normal subgroup of
\(G\), there is a unique \(A\)-group structure on \(F\backslash G\)
for which \(p:G\to F\backslash G\) is a morphism of \(A\)-groups.}
\end{enumerate}

Assertions \SGARefLink{thm:sga3-VIA-3.3.2:item:i}{(i)} and \SGARefLink{thm:sga3-VIA-3.3.2:item:ii}{(ii)} follow from \SGARefLink{thm:sga3-VIA-3.2}{3.2} and \SGARefLink{thm:sga3-VIA-3.3.1}{3.3.1}. Since
\(G\to F\backslash G\) is faithfully flat and locally of finite
presentation, assertion (iii) follows from EGA~IV, 2.2.14, 2.7.1,
and 17.7.7. Assertions (iv) and (v) will be proved in \hyperref[sec:sga3-VIA-5]{Section~5}.
Let us record at once the following corollary.

\medskip
\phantomsection\label{cor:sga3-VIA-3.3.3}%
\noindent\textbf{Corollary 3.3.3.}
\textit{Let \(A\) be an Artinian local ring, let \(G\) be an
\(A\)-group locally of finite type, and let \(H\) be a closed
subgroup of \(G\), flat over \(A\). Write \(p\) for the morphism
\(G\to G/H\), and write \(\lambda\) (respectively,
\(\operatorname{pr}_1\)) for the morphism \(G\times H\to G\) defined
by \(\lambda(g,h)=gh\) (respectively, the projection
\(G\times H\to G\)). Then, for every open subset \(U\) of \(G/H\),
\[
  \mathcal O(U)
  =
  \bigl\{
    \varphi\in\mathcal O\bigl(p^{-1}(U)\bigr)
    \,\bigm|\,
    \varphi\circ\lambda=\varphi\circ\operatorname{pr}_1
  \bigr\}.
\]
In other words, \(\mathcal O(U)\) is the set of
\(\varphi\in\mathcal O(p^{-1}(U))\) such that
\(\varphi(gh)=\varphi(g)\) for every \(A\)-scheme \(S\) and every
\(g\in G(S)\), \(h\in H(S)\).}

Indeed, \(p:G\to G/H\) is faithfully flat and locally of finite
presentation, hence a covering for the \((\mathrm{fppf})\) topology;
the result follows from \SGARefLink{sga3:iv:remark:3:3:3:2}{Expos\'e~IV, 3.3.3.2}.
